\pagestyle{fancy} 
\section{Conclusion}

This thesis asked how ephemeral, task-scoped generative UI components can be designed to improve existing applications without disrupting the user experience. To address that question, it developed a conceptual framework for ephemeral support, implemented the framework in an experimental web artifact, and evaluated the artifact against a persistent baseline in a within-subjects user study. The evaluation provides a qualified answer to that question: ephemeral support can be introduced in a bounded way that preserves users' sense of control, but the implementation tested here did not produce clear performance gains without also introducing additional intrusiveness.

Across the study measures, the ephemeral condition showed a slightly higher task completion rate, although that difference was not statistically significant, and a significantly higher perceived helpfulness rating. Completion time likewise did not improve. By contrast, interaction count increased substantially and perceived intrusiveness increased significantly, while perceived control remained stable and overall preference did not clearly favour either condition. These findings suggest that the main contribution of the ephemeral layer in this prototype was not faster task execution, but a clearer sense of perceived support that came with a noticeable interaction cost.

The thesis's individual contribution lies not only in evaluating ephemeral support as an idea, but in operationalising it through a constrained component-specification model, a validation pipeline for generated interface specifications, a bounded ephemeral overlay approach that preserves host-interface continuity, and an interactive survey artifact that makes empirical evaluation possible.

The broader implication is that ephemeral, task-scoped generative UI remains a promising design direction, but its value depends on more than simply making support temporary. For such support to improve existing applications in practice, it must appear at the right moment, be immediately legible, and offer enough value that users choose to engage with it rather than dismiss it. The present findings support a narrower conclusion than claiming a best-fit domain such as productivity work in general or technical design work specifically. The safer conclusion is that ephemeral support is most credible for bounded, low-consequence subtasks within stable interfaces, particularly where lightweight contextual guidance may help and the cost of interruption remains manageable. The framework proposed in this thesis helps define those boundaries, while the study shows both their promise and their current limits. Technical iterative-refinement settings, such as design exploration or parameter tuning with immediate feedback, are plausible next contexts for evaluation, but they remain future directions rather than validated conclusions from the present study. Future work should therefore test richer scenarios, broader user populations, and more refined support behaviours in order to determine when ephemeral generative assistance becomes not only possible within stable interfaces, but meaningfully beneficial.
